\section{详细设计与实现}
\subsection{空间放置与副本管理}
映射器提供 First Fit、Best Fit 和 aspect-aware 三类装箱策略。基线按顺序选择首个可用位置；优化方案以 Best Fit 为主体，通过剩余空间、形状碎片、分组一致性和已有负载选择候选位置。对于热点专家，副本数由频率、争用、burstiness 和转移影响形成的压力得分决定，并受到额外物理体积预算限制。副本优先分散到不同 Sub-Cube，以避免将逻辑副本重新聚集为同一资源瓶颈。

MoE 优化首先进行多起点互斥分组，再按组需求和组间共现构造分配问题。热点组可在 top-$k$ 范围内用分支定界精确分配，剩余组采用局部迁移、交换和模拟退火改进。该过程并不替代物理装箱：组分配只提供放置偏好，最终仍由 3D 坐标可行性决定是否成功放置。

\subsection{静态调度与约束验证}
模拟器将每个推理实例展开为任务图。仅被轨迹激活的路由专家生成计算任务，共享专家并入激活集合。就绪任务按照 FIFO 或关键路径优先策略选择候选物理副本；调度器计算每个候选副本的依赖释放时间、Sub-Cube 空闲时间、切换惩罚、z 层代价和带宽压力，并选择最早可完成的方案。可选的线性或带宽/z 感知非线性模型用于估计传输/计算重叠收益。

独立验证器不依赖模拟器内部状态，检查 Weight-Cube 坐标边界、三维空间重叠、物理副本 ID 唯一性、schedule 与 placement 的 Sub-Cube 一致性、时间区间互斥、零长度任务、未放置 Cube 和映射率。运行清单记录输入 SHA-256、配置、运行环境与结果摘要，形成可复现证据链。
